% GENERATED FILE. Edit canonical lessons, then run npm run generate:books.
% canonical-source-hash: fnv1a64:8ba718f15d30bf9c
% canonical-lessons: GE-C457-duene, GE-C457-ebbe, GE-C457-flut, GE-C457-bucht, GE-C457-muschel

\chapter{An der Küste}
\label{ch:ge-an-der-kuste}
\hlchaptermodality{457}

\begin{chapteropening}
\textbf{By the end of this chapter:} \emph{I can describe the coast: the dune, low tide, high tide, the bay and the shell.}

Complete the last lesson of chapter 457: I can describe the coast: the dune, low tide, high tide, the bay and the shell.
\end{chapteropening}

\section[die Düne]{die Düne — the dune}

\label{lesson:GE-C457-duene}



\begin{quote}
Before the new one: say the German for the moss, then the German for to rustle.
\end{quote}

\subsection*{You'll want to know: die Düne}
\textbf{die Düne} — \textquotedblleft{}the dune\textquotedblright{}.

\begin{quote}
Hinter den Dünen liegt der Strand.
\end{quote}

\begin{quote}
\emph{The beach lies behind the dunes.}
\end{quote}

Plural: \emph{die Dünen}.

\begin{cousinweb}
\emph{Düne} came into the standard language from Low German, a cousin of Dutch \emph{duin} (EtymDB lists the word); English \emph{dune} came from the same Dutch source.
\end{cousinweb}

\begin{grammarlens}[title={What you've built}]
One more word for telling, explaining and describing at B1.
\end{grammarlens}

\subsection*{Your turn}
\begin{itemize}
\raggedright
  \item \emph{Say it:} \emph{die Düne}
  \item \emph{Say it:} \emph{die Düne}, once more
  \item \emph{Say it:} a sentence of your own with \emph{die Düne}
  \item \emph{Recall:} say the German for the moss, then the German for to rustle, then say \emph{die Düne} again
\end{itemize}

\begin{tcolorbox}[breakable,colback=teal!4,colframe=teal!35!black,title={Before you move on}]
What does die Düne mean? (\textquotedblleft{}The dune\textquotedblright{}.) Say it once more.
\end{tcolorbox}
\section[die Ebbe]{die Ebbe — the low tide, the ebb}

\label{lesson:GE-C457-ebbe}



\begin{quote}
Before the new one: say the German for to rustle, then the German for the dune.
\end{quote}

\subsection*{You'll want to know: die Ebbe}
\textbf{die Ebbe} — \textquotedblleft{}the low tide, the ebb\textquotedblright{}.

\begin{quote}
Bei Ebbe kann man weit ins Watt laufen.
\end{quote}

\begin{quote}
\emph{At low tide you can walk a long way out onto the mudflats.}
\end{quote}

Plural: \emph{die Ebben}, rare. Figurative: \emph{In meinem Portemonnaie ist Ebbe} — my purse is empty.

\begin{cousinweb}
\emph{Ebbe} came into the standard language from Middle Low German \emph{ebbe} (EtymDB), a cousin of English \emph{ebb}.
\end{cousinweb}

\begin{grammarlens}[title={What you've built}]
One more word for telling, explaining and describing at B1.
\end{grammarlens}

\subsection*{Your turn}
\begin{itemize}
\raggedright
  \item \emph{Say it:} \emph{die Ebbe}
  \item \emph{Say it:} \emph{die Ebbe}, once more
  \item \emph{Say it:} a sentence of your own with \emph{die Ebbe}
  \item \emph{Recall:} say the German for to rustle, then the German for the dune, then say \emph{die Ebbe} again
\end{itemize}

\begin{tcolorbox}[breakable,colback=teal!4,colframe=teal!35!black,title={Before you move on}]
What does die Ebbe mean? (\textquotedblleft{}The low tide, the ebb\textquotedblright{}.) Say it once more.
\end{tcolorbox}
\section[die Flut]{die Flut — the high tide; the flood}

\label{lesson:GE-C457-flut}



\begin{quote}
Before the new one: say the German for the dune, then the German for the low tide.
\end{quote}

\subsection*{You'll want to know: die Flut}
\textbf{die Flut} — \textquotedblleft{}the high tide; the flood\textquotedblright{}.

\begin{quote}
Wenn die Flut kommt, ist der Strand ganz schmal.
\end{quote}

\begin{quote}
\emph{When the tide comes in, the beach is very narrow.}
\end{quote}

Plural: \emph{die Fluten}. \emph{Ebbe und Flut} — the tides. Figurative: \emph{eine Flut von E-Mails} — a flood of emails.

\begin{cousinweb}
\emph{Flut} is inherited from Old High German \emph{fluot} and Proto-Germanic \emph{flōduz} (EtymDB); English \emph{flood} is a cousin.
\end{cousinweb}

\begin{grammarlens}[title={What you've built}]
One more word for telling, explaining and describing at B1.
\end{grammarlens}

\subsection*{Your turn}
\begin{itemize}
\raggedright
  \item \emph{Say it:} \emph{die Flut}
  \item \emph{Say it:} \emph{die Flut}, once more
  \item \emph{Say it:} a sentence of your own with \emph{die Flut}
  \item \emph{Recall:} say the German for the dune, then the German for the low tide, then say \emph{die Flut} again
\end{itemize}

\begin{tcolorbox}[breakable,colback=teal!4,colframe=teal!35!black,title={Before you move on}]
What does die Flut mean? (\textquotedblleft{}The high tide; the flood\textquotedblright{}.) Say it once more.
\end{tcolorbox}
\section[die Bucht]{die Bucht — the bay, the cove}

\label{lesson:GE-C457-bucht}



\begin{quote}
Before the new one: say the German for the low tide, then the German for the high tide.
\end{quote}

\subsection*{You'll want to know: die Bucht}
\textbf{die Bucht} — \textquotedblleft{}the bay, the cove\textquotedblright{}.

\begin{quote}
In der kleinen Bucht ist das Wasser ruhig und warm.
\end{quote}

\begin{quote}
\emph{In the little bay the water is calm and warm.}
\end{quote}

Plural: \emph{die Buchten}.

\begin{cousinweb}
\emph{Bucht} came into the standard language from Low German \emph{bucht} 'bend, curve', related to \emph{biegen} 'to bend' (EtymDB lists Dutch \emph{bocht} 'bend' as a cousin): the coast bends inward.
\end{cousinweb}

\begin{grammarlens}[title={What you've built}]
One more word for telling, explaining and describing at B1.
\end{grammarlens}

\subsection*{Your turn}
\begin{itemize}
\raggedright
  \item \emph{Say it:} \emph{die Bucht}
  \item \emph{Say it:} \emph{die Bucht}, once more
  \item \emph{Say it:} a sentence of your own with \emph{die Bucht}
  \item \emph{Recall:} say the German for the low tide, then the German for the high tide, then say \emph{die Bucht} again
\end{itemize}

\begin{tcolorbox}[breakable,colback=teal!4,colframe=teal!35!black,title={Before you move on}]
What does die Bucht mean? (\textquotedblleft{}The bay, the cove\textquotedblright{}.) Say it once more.
\end{tcolorbox}
\section[die Muschel]{die Muschel — the shell, the seashell; the mussel}

\label{lesson:GE-C457-muschel}



\begin{quote}
Before the new one: say the German for the high tide, then the German for the bay.
\end{quote}

\subsection*{You'll want to know: die Muschel}
\textbf{die Muschel} — \textquotedblleft{}the shell, the seashell; the mussel\textquotedblright{}.

\begin{quote}
Am Strand sammeln die Kinder Muscheln.
\end{quote}

\begin{quote}
\emph{The children collect shells on the beach.}
\end{quote}

Plural: \emph{die Muscheln}. Also the mussel you eat.

\begin{cousinweb}
\emph{Muschel} is an early borrowing, Old High German \emph{muscula}, from Latin \emph{mūsculus} 'little mouse, muscle, mussel', from \emph{mūs} 'mouse' (EtymDB; Lewis \& Short).
\end{cousinweb}

\begin{grammarlens}[title={What you've built}]
One more word for telling, explaining and describing at B1.
\end{grammarlens}

\subsection*{Your turn}
\begin{itemize}
\raggedright
  \item \emph{Say it:} \emph{die Muschel}
  \item \emph{Say it:} \emph{die Muschel}, once more
  \item \emph{Say it:} a sentence of your own with \emph{die Muschel}
  \item \emph{Recall:} say the German for the high tide, then the German for the bay, then say \emph{die Muschel} again
\end{itemize}

\begin{tcolorbox}[breakable,colback=teal!4,colframe=teal!35!black,title={Before you move on}]
What does die Muschel mean? (\textquotedblleft{}The shell, the seashell; the mussel\textquotedblright{}.) Say it once more.
\end{tcolorbox}
